\section{Proofs}

\noindent
\textbf{Theorem~\ref{thm:rulerunner-correct}.}
Let $\theta$ be a formula in $\F_{RR}$ and $\boldsymbol{\sigma}$ a non-empty
finite trace. Let $RR_\theta$ be the deterministic finite-state monitor induced
by RuleRunner rules. Then $RR_\theta$ returns \textsc{satisfy} iff
$\boldsymbol{\sigma}\models\theta$. Moreover, $RR_\theta$ is prefix-sound: if it
outputs \textsc{satisfy} (resp.\ \textsc{violate}) after reading the prefix
$\sigma_0\dots\sigma_t$ of $\boldsymbol{\sigma}$, then
$\boldsymbol{\sigma}'\models\theta$ (resp.\ $\boldsymbol{\sigma}'\not\models\theta$)
for every trace $\boldsymbol{\sigma}'$ sharing the same prefix.

\begin{proof}
In the following, given a time step $t$, we write $v_t(\theta) \in \{T,F,?\}$ for the truth value assigned to the $[\theta]$ by the evaluation fixpoint of time step $t$, ignoring the mode qualifier.
%
We first show that, if the initial activation of a propositional formula $\pi$ is in the state at the beginning of time step $t$, then the evaluation fixpoint of that cell assigns to every subformula of $\pi$ exactly one truth value, and $v_t(\pi) = T$ iff $\boldsymbol{\sigma},t\models\pi$. In particular, a propositional formula is never undecided. The proof is by structural induction on $\pi$.

For $\pi = \top$ (resp.\ $\pi = \bot$) the single rule $R[\top]\rightarrow[\top]T$ (resp.\ $R[\bot]\rightarrow[\bot]F$) fires. For $\pi = p$, the observation $p$ belongs to the time step iff $p \in \sigma_t$, so exactly one of $R[p],p \rightarrow [p]T$ and
$R[p],\lnot p \rightarrow [p]F$ fires. For $\pi = \lnot\pi'$ the negation
rule is the dual of the rules for $\pi'$. Finally, for $\pi = \pi_1\land\pi_2$, both
arguments are definite by inductive hypothesis, hence the truth table is exactly that of classical conjunction; the rows consuming a $?$ are not enabled, and the modes $L$ and $R$ are
unreachable, because $R[\pi]L$ and $R[\pi]R$ are produced only by the reactivation of $[\pi]?L$ and $[\pi]?R$.

By this argument, we know that a propositional formula which is instantiated at time step $t$ is always evaluated to a definite value within the same time step $t$, and it is discarded, together with every other truth register, at the end of that time step. Note that $X$, $W$, and $U$ reinstalls only a propositional formula. Hence all the instances of a propositional subformula that
are live at time step $t$, possibly created by different temporal operators, refer
to the same time point $t$ and agree on the value they write in the shared
register. This is exactly the property that fails in
Section~\ref{sec:rulerunner-limitation}.

Now, we now prove the following statement by structural induction on
$\theta\in\F_{RR}$. Let $\boldsymbol{\sigma}$ be a non-empty trace and consider
the run of the RuleRunner system of $\theta$ started from its initial activation
at time step $0$. Then:
\begin{enumerate}
    \item[(i)] at every cell $t$ at which the node of $\theta$ is active, the register $[\theta]$ receives exactly one value $v_t(\theta)$;
    \item[(ii)] if $v_t(\theta)\neq{?}$, then $v_t(\theta)=T$ iff $\boldsymbol{\sigma},0\models\theta$;
    \item[(iii)] if $v_t(\theta)={?}$ and $t<last$, the node of $\theta$ is active again at cell $t+1$ with its mode recorded, and so are the arguments it has to read there;
    \item[(iv)] if $v_t(\theta)={?}$ for every $t\leq last$, then the END resolution of $[\theta]$ returns true iff $\boldsymbol{\sigma},0\models\theta$.
\end{enumerate}

\noindent
\textbf{(Base case.)}

\begin{itemize}
    \item $\theta=\pi$. By the propositional case, $v_0(\pi)$ is definite and $v_0(\pi)=T$ iff $\boldsymbol{\sigma},0\models\pi$, which gives (i) and (ii); (iii) and (iv) are vacuous, and the node is not active after cell $0$.
    \item $\theta=X\pi$. The rule $R[X\pi]\rightarrow[X\pi]?$ fires and no other rule of the node is enabled, so $v_0(X\pi)={?}$ and (i) holds, while (ii) is vacuous. If $0=last$, then $\boldsymbol{\sigma},0\not\models X\pi$, and the END resolution returns false, which gives (iv). If $0<last$, the reactivation of the unqualified $[X\pi]?$ installs $R[X\pi]M$ together with the initial activation of $\pi$, which gives (iii). By the propositional case $v_1(\pi)=T$ iff $\boldsymbol{\sigma},1\models\pi$, and $R[X\pi]M,[\pi]V\rightarrow[X\pi]V$, so $v_1(X\pi)$ is definite and it is true iff $\boldsymbol{\sigma},0\models X\pi$: this makes (iv) vacuous.
    \item $\theta=W\pi$. Same as the previous case, except that the END resolution of the unqualified $[W\pi]?$ returns true, which is correct.
    \item $\theta=\pi_1\muntil\pi_2$. The node is activated at time step $0$ in mode $A$, together with the initial activations of both arguments, and its reactivation re-installs both of them at every time step. Since by the propositional case the arguments are always definite, $?$ is never returned as truth value and the modes $B$, $L$ and $R$ are unreachable for the whole run. Hence, we can have
    $v_t(\pi_1\muntil\pi_2)$ equals true if $\boldsymbol{\sigma},t\models\pi_2$, false if $\boldsymbol{\sigma},t\not\models\pi_2$ and $\boldsymbol{\sigma},t\not\models\pi_1$, and $?A$ otherwise. This gives (i) and (iii).

    We prove by induction on $t$ the invariant: \emph{the node is active at cell $t$ iff $\boldsymbol{\sigma},k\models\pi_1$ and $\boldsymbol{\sigma},k\not\models\pi_2$ for every $k<t$}. It holds trivially at $t=0$, and it is preserved because the node is active at cell $t+1$ only if $v_t={?A}$, i.e.\ only if $\pi_1$ holds and $\pi_2$ fails at $t$. Under the invariant, (ii) follows: if $v_t=T$ then $\boldsymbol{\sigma},t\models\pi_2$ and $\pi_1$ holds at every $k<t$, so $t$ witnesses $\boldsymbol{\sigma},0\models\pi_1\muntil\pi_2$; if $v_t=F$ then no $j<t$ witnesses the until by the invariant, $j=t$ fails because $\pi_2$ does, and no $j>t$ can witness it because $\pi_1$ fails at $t$, hence $\boldsymbol{\sigma},0\not\models\pi_1\muntil\pi_2$. Finally, (iv): if the node is undecided at every cell, then by the invariant $\pi_2$ fails everywhere up to $last$, so $\boldsymbol{\sigma},0\not\models\pi_1\muntil\pi_2$; and the END resolution of an until returns the END value of its right argument, which is the definite $[\pi_2]F$ recorded at cell $last$.
\end{itemize}

\noindent
\textbf{(Inductive case.)}

\begin{itemize}
    \item $\theta=\lnot\theta'$. The negation rules map $T$ to $F$, $F$ to $T$ and $?$ to $?$, so (i) holds and $v_t(\theta)\neq{?}$ iff $v_t(\theta')\neq{?}$, in which case (ii) follows from the inductive hypothesis on $\theta'$. If $v_t(\theta')={?}$, then $[\lnot\theta']?\rightarrow R[\lnot\theta']$ carries the wrapper while $\theta'$ reactivates itself through its own rules, which together with the inductive hypothesis (iii) gives (iii). Likewise, the END resolution of $\lnot\theta'$ is the complement of the END resolution of $\theta'$, so (iv) follows from the inductive hypothesis (iv).
    \item $\theta=\theta_1\land\theta_2$. Consider the first cell $t$ at which one of the two arguments becomes definite; before $t$ both are undecided, the row $(?,?)$ of the mode-$B$ table returns $?B$, and both arguments are carried and reactivate through their own rules by the inductive hypothesis (iii). If no such cell exists, (iv) holds because the END resolution in mode $B$ recurses on both arguments, which are correct by the inductive hypothesis (iv), and the conjunction of their END values is the value of $\theta$ at time $0$. Otherwise, we distinguish three cases at cell $t$.
    If both arguments are definite, the mode-$B$ table returns their classical conjunction, and (ii) follows from the inductive hypothesis (ii) applied to both.
    If exactly one is definite and its value is $F$ --- say $v_t(\theta_1)=F$ --- the row $(F,?)$ returns $F$, and by the inductive hypothesis (ii) $\boldsymbol{\sigma},0\not\models\theta_1$, hence $\boldsymbol{\sigma},0\not\models\theta$ and (ii) holds.
    If exactly one is definite and its value is $T$ --- say $v_t(\theta_1)=T$ --- the row $(T,?)$ returns $?R$, so from cell $t+1$ on only $R[\theta]R$ is carried and $\theta_1$ is dropped from monitoring; dropping it is correct, because by the inductive hypothesis (ii) $\boldsymbol{\sigma},0\models\theta_1$ permanently, and $\theta_1$'s own subsystem is not reactivated, its register being definite. The mode-$R$ rules replicate the mode-$B$ column at the pin value $T$ of the left argument, so $v_{t'}(\theta)=v_{t'}(\theta_2)$ at every cell $t'\geq t$, and (ii), (iii) follow from the inductive hypothesis on $\theta_2$ together with $\boldsymbol{\sigma},0\models\theta$ iff $\boldsymbol{\sigma},0\models\theta_2$. For (iv), the END resolution in mode $R$ pins the left argument to $T$, which is the very value it was assigned at cell $t$, and recurses on $\theta_2$, so the claim follows from the inductive hypothesis (iv). The case in which $\theta_2$ is the argument that settles is symmetric, with $?L$ in place of $?R$ and the right argument pinned to $T$.
\end{itemize}

\noindent
In summary, if the value of $\theta$ becomes definite at some time step $t\leq last$, the monitor
halts there and, by (ii), its verdict agrees with
$\boldsymbol{\sigma}\models\theta$ (this ensures also prefix soundness).
Otherwise the register is undecided at every cell, and by (iv) the END
resolution returns the correct value. In both cases $RR_\theta$ returns
\textsc{satisfy} iff $\boldsymbol{\sigma}\models\theta$.
\end{proof}


\noindent
\textbf{Corollary~\ref{cor:rulerunner-certifier}.}
If $\varphi \in \F_{RR}$, then $\varphi \in \C_{RR}$.

\begin{proof}
It follows directly by Theorem \ref{thm:rulerunner-correct}.
\end{proof}


\noindent
\textbf{Lemma~\ref{lem:horizon-locality}.}
Let $\varphi$ be a formula with $h(\varphi)=h<\infty$. Let $l=\min(t+h,last)$ and $\boldsymbol{w}=\boldsymbol{\sigma}[t..l]$. For every non-empty trace $\boldsymbol{\sigma}$ and every position $t\leq last$,
\[
\boldsymbol{\sigma},t\models\varphi
\,\text{ if and only if } \,
\boldsymbol{w},0\models\varphi
\]

\begin{proof}
The proof is by structural induction on $\varphi$.

If $\varphi$ is $\top$, $\bot$ or an atom, the statement is trivially satisfied.

If $\varphi$ is $\lnot\varphi'$, we have
$h'=h(\varphi')\leq h$. It is easy to see that $\min(t+h',last)=\min(t+h',l)=t+\min(h',l-t)$. 
The inductive hypothesis applied to $\boldsymbol{\sigma}$
at $t$ reduces $\varphi'$ to the window $\boldsymbol{\sigma}[t,\min(t+h',last)]$, and applied to $\boldsymbol{w}$ at $0$ it reduces it to $\boldsymbol{w}[0,\min(h',l-t)]$.
Since $\varphi'$ has the same value in $\boldsymbol{\sigma}$ at $t$ and in
$\boldsymbol{w}$ at $0$, so does $\varphi$.
%
If $\varphi$ is $\varphi_1\circ\varphi_2$, a similar argument as above applies.

If $\varphi$ is $X\varphi'$, we have $h'=h(\varphi')$ so that $h=1+h'$. When $t=last$ we
have $l=t$ and $\boldsymbol{w}$ has a single position, so both
$\boldsymbol{\sigma},t\models X\varphi'$ and
$\boldsymbol{w},0\models X\varphi'$ are false. When $t<last$ we have
$l\geq t+1$, so both traces have a successor position and
the two sides amount to $\boldsymbol{\sigma},t+1\models\varphi'$ and
$\boldsymbol{w},1\models\varphi'$. The inductive hypothesis reduces them to $\boldsymbol{\sigma}[t+1,\min(t+1+h',last)]$ and
$\boldsymbol{w}[1,\min(1+h',l-t)]$, which is the same trace. The case $\varphi=W\varphi'$ is
the same, except that a single-position $\boldsymbol{w}$ makes both sides true.
\end{proof}

\noindent
\textbf{Theorem~\ref{thm:rulerunner-bounded-correct}.}
Let $\alpha = sk(\varphi)$. If $RR_\alpha$ is language-equivalent to $A_\alpha$,
then for every finite trace $\boldsymbol{\sigma}$ the composition of the
bounded-event transducer and $RR_\alpha$ returns \textsc{satisfy} iff
$\boldsymbol{\sigma}\models\varphi$. Moreover, if $RR_\alpha$ is prefix-sound,
then so is the composition.

\begin{proof}
The empty trace is settled by the separately computed boundary value, since it
has no position at which an event could be emitted, so let
$\boldsymbol{\sigma}$ be non-empty. We prove the claim in two steps.

\medskip
\noindent
\textbf{(The transducer emits the derived trace.)} Consider a position $p$
emitted while the trace is being read, so that its buffer is
$\boldsymbol{\sigma}[p..p+H]$ with $p+H\leq last$. For every island
$\varphi_i$, Lemma~\ref{lem:horizon-locality} applied to $\boldsymbol{\sigma}$
at $p$ reduces $\boldsymbol{\sigma},p\models\varphi_i$ to the window
$\boldsymbol{\sigma}[p..p+h(\varphi_i)]$, and applied to the buffer at $0$ it
reduces the transducer's own evaluation to the same window, because
$h(\varphi_i)\leq H$. The two values therefore coincide. Consider now a
position $p$ emitted during the flush: its buffer is exactly the suffix
$\boldsymbol{\sigma}[p..last]$, and evaluating a formula at position $0$ of a
suffix is evaluating it at $p$ in $\boldsymbol{\sigma}$, the truncated
readings of $X$ and $W$ being the end-of-trace semantics of \ltlf. In both
cases the emitted cell also carries $\sigma_p$ unchanged, and the cells are
emitted in increasing order of $p$, one per input cell. Hence the transducer
outputs exactly $\widehat{\boldsymbol{\sigma}}$.

\medskip
\noindent
\textbf{(Substitution.)} We show that $\boldsymbol{\sigma},t\models\varphi$ iff
$\widehat{\boldsymbol{\sigma}},t\models\alpha$ for every position $t$, by
structural induction on $\varphi$ following the rewriting that produces
$\alpha$. If the node is one of the islands $\varphi_i$, it is replaced by
$e_i$, and $e_i\in\hat{\sigma}_t$ iff $\boldsymbol{\sigma},t\models\varphi_i$
by the definition of the derived trace. If the node is an atom, it is left
unchanged, and $\hat{\sigma}_t\cap\calP=\sigma_t$. Otherwise the node is
unbounded, the rewriting preserves its operator and recurses on its arguments,
and the inductive hypothesis applies to them; the positions at which they are
evaluated are the same in the two traces, because
$\widehat{\boldsymbol{\sigma}}$ and $\boldsymbol{\sigma}$ have the same length,
hence the same $last$.

\medskip
\noindent
Taking $t=0$, the skeleton monitor consumes $\widehat{\boldsymbol{\sigma}}$
and, by language equivalence, returns \textsc{satisfy} iff
$\widehat{\boldsymbol{\sigma}}\models\alpha$, that is iff
$\boldsymbol{\sigma}\models\varphi$.

Assume now that $RR_\alpha$ is also prefix-sound, and suppose that it emits a
definite verdict after the transducer has read $\sigma_0,\ldots,\sigma_t$. At
that point the transducer has emitted the prefix
$\hat{\sigma}_0,\ldots,\hat{\sigma}_{t-H}$ of the derived trace, which is empty
when $t<H$. Prefix soundness makes that verdict correct for every trace over
$\calP\cup\{e_1,\ldots,e_m\}$ extending the emitted prefix. Now, every
non-empty trace $\boldsymbol{\sigma}'$ extending $\sigma_0,\ldots,\sigma_t$
derives a trace $\widehat{\boldsymbol{\sigma}'}$ that extends that prefix: each
emitted position $p\leq t-H$ has its whole window $\boldsymbol{\sigma}[p..p+H]$
inside the part already read, so by Lemma~\ref{lem:horizon-locality} it
receives the same value in $\boldsymbol{\sigma}'$, whether the extension is
proper or $\boldsymbol{\sigma}'$ stops at $t$ and $p$ is flushed. The
realizable derived continuations are therefore a subset of the continuations
quantified over by prefix soundness, and by substitution the verdict agrees
with $\boldsymbol{\sigma}'\models\varphi$ for all of them.
\end{proof}

\noindent
\textbf{Corollary~\ref{cor:rulerunner-bounded-fragment}.}
For every $\varphi\in\F_{BE}$, the bounded-event RuleRunner returns the correct
final verdict, and it is prefix-sound.

\begin{proof}
By definition, if $\varphi\in\F_{BE}$ its skeleton $sk(\varphi)$ belongs to
$\F_{RR}$. By Corollary~\ref{cor:rulerunner-certifier}, every formula in
$\F_{RR}$ is in $\C_{RR}$, so $RR_{sk(\varphi)}$ is language-equivalent to
$A_{sk(\varphi)}$; and by Theorem~\ref{thm:rulerunner-correct} it is also
prefix-sound. Thus, we
can conclude that the bounded-event RuleRunner returns \textsc{satisfy} iff
$\boldsymbol{\sigma}\models\varphi$ and it is also prefix-sound.
\end{proof}


\noindent
\textbf{Theorem~\ref{thm:rulerunner-bounded-online}.}
After every non-empty prefix, the extrapolated bounded-event monitor returns a
three-valued label identical to that of the canonical DFA for the original
formula.

\begin{proof}
By Theorem~\ref{thm:rulerunner-bounded-correct},
$\calL(P_\varphi)=\calL(A_\varphi)$. In a deterministic automaton, the state
reached after a prefix $\boldsymbol{u}$ is an accepting sink if and only if
every continuation of $\boldsymbol{u}$ belongs to the recognised language, and a
trap if and only if no continuation does; both conditions depend only on
$\boldsymbol{u}$ and on the language, not on the automaton recognising it. Since
$P_\varphi$ and $A_\varphi$ recognise the same language and read the same
alphabet, the state reached by $P_\varphi$ after $\boldsymbol{u}$ is an
accepting sink, respectively a trap, exactly when the state reached by
$A_\varphi$ is. The extrapolated monitor returns the label of the former and the
symbolic monitor of Section~\ref{sec:symbolic} returns the label of the latter,
so the two labels coincide.
\end{proof}



\noindent
\textbf{Theorem~\ref{thm:rulerunner-progression-correct}.}
Let $\varphi$ be an \ltlf formula and let $RR^{\mathsf{prog}}_\varphi$ be the (eager) progression-based RuleRunner monitor constructed as above. Then, for every finite  trace $\boldsymbol{\sigma}$, the monitor returns \textsc{satisfy} iff $\boldsymbol{\sigma}\models\varphi$, and it returns \textsc{violate} iff $\boldsymbol{\sigma}\not\models\varphi$. Moreover, any verdict returned before the last cell is sound for all finite continuations of the prefix read so far.

\begin{proof}
Let $\boldsymbol{\sigma}=(\sigma_0,\ldots,\sigma_{T-1})$ be a finite trace. For $0\leq t\leq T$, let $\Sigma_t$ be the active root set after the prefix $\sigma_0\cdots\sigma_{t-1}$ has been consumed (and thus before reading $\sigma_t$ when $t<T$). Let
\[
    \rho_t = \bigwedge_{\chi\in\Sigma_t}\chi
\]
be the residual it denotes. By construction,
\[
    \Sigma_0=\mathsf{split}_\wedge(\mathsf{nf}(\varphi)),
\]
and therefore $\rho_0$ is equivalent to $\varphi$.

We first use the defining property of progression. For every formula $\chi$, every cell $s$, and every finite continuation $\boldsymbol{\tau}$, including the empty continuation $\epsilon$,
\begin{equation}
\label{eq:progression-step}
    s\cdot\boldsymbol{\tau}\models\chi
    \quad\Longleftrightarrow\quad
    \boldsymbol{\tau}\models\mathsf{prog}(\chi,s).
\end{equation}
This property follows by structural induction on $\chi$. The atomic and Boolean cases are immediate from the definition of $\mathsf{prog}$. For $X\chi$, $s\cdot\boldsymbol{\tau}\models X\chi$ iff $\boldsymbol{\tau}$ is non-empty and $\boldsymbol{\tau}\models\chi$. Since $F\top$ holds exactly on non-empty finite traces, this is equivalent to $\boldsymbol{\tau}\models\chi\land F\top$. Dually, $W\chi$ is true either when the suffix is empty or when $\chi$ holds at its first cell; $G\bot$ holds exactly on the empty trace, yielding $\mathsf{prog}(W\chi,s)=\chi\lor G\bot$. For until, $s\cdot\boldsymbol{\tau}\models\chi_1U\chi_2$ iff either $\chi_2$ holds at the current cell, or $\chi_1$ holds at the current cell and $\chi_1U\chi_2$ holds on the suffix; by the induction hypothesis this is equivalent to
\[
    \boldsymbol{\tau}\models
    \mathsf{prog}(\chi_2,s)\lor
    \bigl(\mathsf{prog}(\chi_1,s)\land(\chi_1U\chi_2)\bigr).
\]
The derived operators inherit the property from their definitions.

We now prove the residual invariant:
\begin{equation}
\label{eq:rulerunner-residual-invariant}
    \sigma_0\cdots\sigma_{t-1}\cdot\boldsymbol{\tau}\models\varphi
    \quad\Longleftrightarrow\quad
    \boldsymbol{\tau}\models\rho_t
\end{equation}
for every $0\leq t\leq T$ and every finite continuation $\boldsymbol{\tau}$.

For $t=0$, the claim holds because $\rho_0$ is equivalent to $\varphi$. Assume it holds for some $t<T$. By the reactivation rule,
\[
    \rho_{t+1}
    \equiv
    \mathsf{nf}\left(
        \bigwedge_{\chi\in\Sigma_t}\mathsf{prog}(\chi,\sigma_t)
    \right).
\]
Since $\rho_t=\bigwedge_{\chi\in\Sigma_t}\chi$ and progression distributes over conjunction up to logical equivalence, this is equivalent to
\[
    \rho_{t+1}
    \equiv
    \mathsf{nf}\bigl(\mathsf{prog}(\rho_t,\sigma_t)\bigr).
\]
Therefore, for every  continuation $\boldsymbol{\tau}$,
\[
\begin{array}{rcl}
\sigma_0\cdots\sigma_t\cdot\boldsymbol{\tau}\models\varphi
&\Longleftrightarrow&
\sigma_t\cdot\boldsymbol{\tau}\models\rho_t
\quad\text{by the induction hypothesis,}\\
&\Longleftrightarrow&
\boldsymbol{\tau}\models\mathsf{prog}(\rho_t,\sigma_t)
\quad\text{by Equation~\eqref{eq:progression-step},}\\
&\Longleftrightarrow&
\boldsymbol{\tau}\models\rho_{t+1}
\quad\text{because $\mathsf{nf}$ preserves semantics.}
\end{array}
\]
Thus the invariant holds at $t+1$.

After consuming all $T$ cells, repeated application of the invariant gives
\[
    \boldsymbol{\sigma}\models\varphi
    \quad\Longleftrightarrow\quad
    \epsilon\models\rho_T.
\]
The final evaluation rules compute precisely whether the residual accepts the empty trace. Hence the monitor returns \textsc{satisfy} exactly on traces satisfying $\varphi$, and \textsc{violate} otherwise.

Finally, the operational register construction realizes the symbolic transition used in the proof. At time $t$, the state $S_t=sub(\rho_t)$ contains every subformula needed to evaluate and progress the active roots. Since the eager monitor allocates registers for all formulas in $sub(cl(\varphi))$, every such support formula has a corresponding register, and the bottom-up computation can be carried out by finitely many rules. Only the active roots in $\Sigma_t$ are reactivated, so support subformulas do not acquire an independent temporal meaning. Therefore the rule-based monitor implements exactly the residual update used in the invariant.

If an early verdict is returned, the current residual is respectively valid or unsatisfiable over all finite continuations. By Equation~\eqref{eq:rulerunner-residual-invariant}, every completion of the prefix read so far then has the same truth value for $\varphi$, which proves soundness of early termination.
\end{proof} 
